%% =====================================================================
%%  B4-T-10-01 -- what B4 contributes to "Value as a Function of Availability"
%%  -------------------------------------------------------------------
%%  LAYER A: APPEND-ONLY. Written once, then left alone. Material found
%%  only here sits in the `unique` slot and is protected by rule R10.
%% =====================================================================
%% @kind: source
%% @id: B4-T-10-01
%% @book: B4
%% @topic: T-10-01
%% @locator: ch. 7, pp. 205--210
%% @status: distilled
%% @unique: yes
%% @integrated: yes
%% @updated: 2026-09-19

\dossier{B4}{T-10-01}{\bookshort{B4}}{ch. 7, pp. 205--210}

\begin{distilled}
  \item Opportunities increase in apparent value as they become less available.
  \item Rarity is used as a proxy for quality, and the proxy can be produced without the underlying condition.
  \item The stake shifts from acquiring a gain to avoiding the loss of an option, and loss carries more weight than the equivalent gain.
  \item Applied to limited editions, closing windows, restricted eligibility and withdrawal of access.
\end{distilled}

\begin{unique}
  \item The loss-shift account: scarcity works because it converts a gain into a prospective loss, which is why the pressure is felt as appraisal rather than as hurry.
  \item The observation that nothing about the object need change --- only the claim about its availability.
\end{unique}
